
\section{Early Attempts at Solving Sokoban}

In the late 1990s and 2000s, the focus of Sokoban research shifted from the difficulty of the problem to its solution. Early attempts at solving Sokoban relied on variations of traditional algorithmic search such as A*. Junghanns et al.\ presented numerous research projects on Sokoban in the late 1990s. In one of these works, they attempted to solve all 90 puzzles of the XSokoban test suite using IDA* \cite{myers_xsokoban, junghanns1997sokoban, junghanns1998sokoban}. The researchers demonstrated a modified IDA* search method named ``Rolling Stone'' that incorporates a Deadlock Table to identify dead-end configurations earlier in the heuristic search-tree expansion in order to prune the search. The Deadlock Table consisted of 22 million entries of frequently occurring board patterns that lead to a dead end. Rolling Stone also introduced a concept called ``Macro Moves'' that collapses a sequence of forced moves into single steps for more efficient memory use during search while preserving solution optimality.

The Rolling Stone algorithm was able to find solutions to 16 of the 90 problems with a search tree shallower than the solution itself. This was made possible by Macro Moves. However, this method comes with significant limitations, as it requires a 22-million-row deadlock pattern database, a piece of hard-coded domain knowledge. The paper also recognizes that Rolling Stone spends 90\% of its execution time updating the lower bounds, which means most of the processing time is spent on optimizing a found solution \cite{junghanns1998sokoban}.

A similar method was tried by Pereira et al.\ \cite{pereira2013finding}, where the researchers employed a pattern database, a pre-computed lookup table that stores the shortest distances for simplified local problems, to be used as a heuristic in the IDA* search. The authors introduced a concept called instance decomposition. Because it is not possible to know which stone is placed on which goal position before the puzzle is solved, instance decomposition assigns each stone to a specific goal before running the search. This constraint acts as an intermediate goal and converts the multi-goal problem into a single-goal problem. During the search, the Pattern Database provides an admissible heuristic to be used for pruning the search tree. This motivates the search for better heuristics and representations.

\section{Reinforcement Learning Approaches}

Pattern-based search methods remained the dominant approach to Sokoban until Google released I2A, an Imagination-Augmented Agent based on reinforcement learning.

\section{I2A and Boxoban}
In 2017, Google researchers introduced a novel architecture \cite{NEURIPS2017_7ca66fc0} for reinforcement learning (RL) that bridges the gap between model-free and model-based approaches. Unlike standard agents that react directly to the environment, Imagination-Augmented Agents (I2As) learn to "imagine" future steps using an internal, potentially imperfect environment model. Instead of predicting an action solely from the current state, the agent first samples several future scenarios and interprets these trajectories to inform its next move.

Classical planning methods, such as Monte Carlo Tree Search (MCTS), expand search branches based on a model's predictions, which typically requires the model to be perfectly accurate. I2As circumvent this limitation by using an encoder to extract meaningful features from simulated trajectories, making the system significantly more robust to model inaccuracies or failures.

The I2A architecture consists of three major components.

First is the Imagination Core(IC). IC is a type of model that is now more commonly known as World Model.\cite{worldmodel} This model is pretrained on trajectories of the environment. The transition model $M$ takes the current observation $s_t$ and a proposed action $a_t$ to predict the subsequent state $s_{t+1}$ and the immediate reward $r_{t+1}$


\[M(s_{t+1},\ r_{t+1}|s_t,\ a_t)\]

Follwing the IC is the Rollout Stategy. To generate an imagined trajectory, the agent must decide not just the first action, but a sequence of future actions. The rollout stategy is typically a simplified, distilled version of the agent's current policy. It performs $n$ look-ahead steps to produce a rollout consisting of imagined states, actions, and rewards. This component is essential because it determines which parts of the future state-space the agent explores during its deliberation phase.

The third component is the Rollout Encoder. Traditional model-based RL often fails if the internal model is inaccurate as the agent might trust a high-reward path that doesn't actually exist. The Rollout Encoder processes the entire sequence of imagined transitions and compresses them into a fixed dimensional vector. The encoder learns to extract features from the imagination. If the IC is inaccurate, the encoder can learn to identify those inconsistencies. 

Finally, the Aggregator is a network that uses the current environment state and the imagined path from the rollout encoder to select an action. By replacing parts of MCTS with pretrained networks, I2A can perform training tasks such as solving Sokoban puzzles.

Google demonstrated the efficacy of the I2A agent using Boxoban, a massive dataset consisting of 9 million procedurally generated Sokoban levels. Sokoban is a challenging benchmark because moves are often irreversible, requiring deep strategic planning.

The performance of I2A was compared against prominent model-free architectures, specifically A3C (Asynchronous Advantage Actor-Critic). While the A3C agent achieved a 60\% solve rate, the I2A agent solved 85\% of the same test set. Furthermore, I2A proved more resource-efficient; because its training procedure is partially independent of the live environment, it required substantially fewer environment interactions to converge.

Beyond the architecture itself, the Boxoban dataset remains a significant contribution to the RL community, providing a standardized and rigorous benchmark for evaluating generalization and planning in artificial agents.



\subsection{Forward-Backward Reinforcement Learning}

Shoham et al. \cite{shoham2021solving} addressed the sparse reward problem in Sokoban by training two complementary networks: a backward-facing model and a forward-facing model. Because reward in Sokoban is only received when all boxes are on goal positions, the backward model is initialized at the goal state and trained on a reversed task by pulling boxes away from goals rather than pushing them towards them. This reversed task provides dense reward early in training and allows the agent to learn a backward value function $V_B$ efficiently. From this function, a backward trajectory $T$ is constructed for each puzzle instance by initializing at the goal state and using a planning algorithm to find the path that maximizes $V_B$.

The forward-facing value function $V_F$ then conditions on both the current state $s$ and instance-specific hint features derived from the backward trajectory $T$. This conditioning partially alleviates the sparsity of the forward reward signal by providing structural information about where boxes should ultimately be placed. The authors note that despite this novel approach, classical IDA* search still provides the strongest overall performance among the methods compared, underscoring the continued relevance of principled heuristic search for Sokoban.


\subsection{Learning Generalized Reactive Policies}

While I2A method is technically impressive, in order to achieve the best performance, the architecture relied heavily on large amount of data in order to not only determine the best action, but also to learn a world model that learns the environment itself. This means that the bottleneck of this architecture is the computation required to train the world model. Groshev et al. \cite{groshev2018learning} proposed a more computational efficient method in which a deep neural network learns a Generalized Reactive Policy (GRP) that directly maps a problem instance and a current state to an action. The DNN is then used as a heuristic function to an A* search. 

To train the DNN, the authors first generated a dataset of successful solution trajectories for a set of Sokoban levels. Each trajectory consists of a sequence of board states and corresponding actions that lead from the initial state to the goal state. These successful trajectories are generated by a Fast-Forward planner \cite{Hoffmann_2001}, which is a domain-independent heuristic search 
planner based on agumented Breath First Search(BFS). The DNN is then trained on this dataset to predict the successful action given the current board state.

The core contribution of this work is the demonstration that a DNN can learn a heuristic function for A* search based on solution trajectories.

Unlike prior approaches that relied heavily on hand-crafted domain knowledge or manually engineered feature representations, their method trains the policy from scratch using only a set of successful execution traces, reducing dependence on problem-specific expertise. For Sokoban, the authors employed a multi-headed convolutional neural network that jointly predicts both action probability and plan length from the current board state. The success of this architecture on spatial planning tasks provides motivation for the use of convolutional processing as a compression baseline in our own approach.


\section{DeepCubeA}

DeepCubeA, proposed by Agostinelli et al. \cite{agostinelli2019solving}, replaces the hand-crafted heuristic in A* with a deep neural network $h_\theta$ trained to estimate the cost-to-go from any state to the goal. The network is trained through a self-supervised procedure called \textit{Autodidactic Iteration}: beginning from the solved (goal) state, random sequences of actions of increasing length are applied to generate scrambled states. At each depth $d$, the network learns to predict $d$ as the cost-to-go from the resulting state. This backward induction scheme produces a heuristic that approximates true distance across the reachable state space without any hand-engineered domain knowledge.

At inference time, DeepCubeA uses A* (WA*), scaling the learned heuristic by a weight $w \geq 1$:

$$f(n) = g(n) + w \cdot h_\theta(n)$$

A weight $w > 1$ biases node expansion toward the goal, sacrificing the optimality guarantee in exchange for significantly faster convergence — a practical trade-off for large state spaces where optimal search is intractable. DeepCubeA demonstrated strong performance on challenging combinatorial puzzles including the Rubik's Cube and the 15-puzzle, establishing it as a general-purpose learned search framework. A subsequent work extending this approach, which uses Answer Set Programming to specify goals to deep neural networks, demonstrated the framework's applicability to Sokoban, successfully solving highly difficult instances involving more than 20 boxes. \cite{agostinelli2021admissible}
